\section*{Chapter \ref{ch:g6:rocks-and-minerals} --- Rocks and Minerals}

\begin{solution}{exo:g6:rocks-and-minerals:1}
Granite: igneous. Sandstone: sedimentary. Marble: metamorphic. Basalt:
igneous. Limestone: sedimentary. Slate: metamorphic.
\end{solution}

\begin{solution}{exo:g6:rocks-and-minerals:2}
Quartz is a \omterm{def:g6:rocks-and-minerals:mineral}{mineral}; granite is a \omterm{def:g3:stones-pebbles-sand:rock}{rock}. A \omterm{def:g6:rocks-and-minerals:mineral}{mineral} is a single natural
substance with its atoms arranged in a regular way; a \omterm{def:g3:stones-pebbles-sand:rock}{rock} is made of one
or several \omterm{def:g6:rocks-and-minerals:mineral}{minerals} stuck together (granite: feldspar, quartz and mica).
\end{solution}

\begin{solution}{exo:g6:rocks-and-minerals:3}
Talc (1), calcite (3), quartz (7), diamond (10). Diamond, the hardest,
scratches all the others.
\end{solution}

\begin{solution}{exo:g6:rocks-and-minerals:4}
``Molten \omterm{def:g3:stones-pebbles-sand:rock}{rock} becomes an \omterm{def:g6:rocks-and-minerals:rock-families}{igneous rock} by cooling. A limestone becomes a
marble by heat and pressure. A granite turns into grains of \omterm{def:g3:stones-pebbles-sand:sand}{sand} by
\omterm{def:g5:water-shapes-the-land:weathering}{weathering}.''
\end{solution}

\begin{solution}{exo:g6:rocks-and-minerals:5}
(a) Basalt. (b) Marble. (c) Sandstone.
\end{solution}

\begin{solution}{exo:g6:rocks-and-minerals:6}
Its \omterm{def:g6:rocks-and-minerals:hardness}{hardness} is more than that of a fingernail (about 2.5) and less than
that of a copper coin (about 3.5). It could be calcite (3).
\end{solution}

\begin{solution}{exo:g6:rocks-and-minerals:7}
Quartz: $21 \div 60 = 0.35$, that is 35~\%. Feldspar: $33 \div 60 = 0.55$,
that is 55~\%. Mica: $6 \div 60 = 0.10$, that is 10~\%. Check:
$35 + 55 + 10 = 100$.
\end{solution}

\begin{solution}{exo:g6:rocks-and-minerals:8}
Granite cools very slowly, deep underground: its grains have a long time
to grow, and become large. Basalt comes from \omterm{def:g4:volcanoes-earthquakes-first-look:eruption}{lava}, which cools quickly at
the surface: its grains have no time to grow and stay too small to see.
\end{solution}

\begin{solution}{exo:g6:rocks-and-minerals:9}
(a) Limestone is sedimentary; marble is metamorphic. (b) Deep underground,
heat and pressure transformed it without melting it: its calcite grew into
shiny crystals that fit together. (c) Both are made of the same \omterm{def:g6:rocks-and-minerals:mineral}{mineral},
calcite, and it is calcite that fizzes under acid.
\end{solution}

\begin{solution}{exo:g6:rocks-and-minerals:10}
(a) Yes: quartz (7) is harder than glass (about 5.5). (b) No: fluorite (4)
is softer than glass. (c) Yes: a steel nail (about 6.5) is harder than
feldspar (6).
\end{solution}

\begin{solution}{exo:g6:rocks-and-minerals:11}
For example: (1) the granite of the \omterm{def:g1:landscapes-around-us:landform}{mountain} (igneous) is broken up by
\omterm{def:g5:water-shapes-the-land:weathering}{weathering}, and its quartz grains are set free; (2) a \omterm{def:g1:landscapes-around-us:water}{river} carries them
off: \omterm{def:g5:water-shapes-the-land:erosion}{erosion}; (3) the grains are laid down in layers at the bottom of the
\omterm{def:g2:water-on-earth:ocean}{sea}: \omterm{def:g5:water-shapes-the-land:deposition}{deposition}; (4) buried under other layers, they are pressed and
cemented into a sandstone: hardening (the \omterm{def:g3:stones-pebbles-sand:rock}{rock} is now sedimentary);
(5) carried deep underground, the sandstone is transformed by heat and
pressure into a \omterm{def:g6:rocks-and-minerals:rock-families}{metamorphic rock}.
\end{solution}

\begin{solution}{exo:g6:rocks-and-minerals:12}
(a) Not scratched by steel (about 6.5) and able to scratch glass: harder
than 6.5. Scratched by topaz (8): softer than 8. Its \omterm{def:g6:rocks-and-minerals:hardness}{hardness} is 7: it is
quartz. (b) Quartz is harder than almost everything that rubs against it
(other grains, \omterm{def:g3:stones-pebbles-sand:pebble}{pebbles}): it is hardly worn away. Calcite (3) and mica
(2 to 3) are much softer: rubbing grinds them into tiny pieces that are
carried away.
\end{solution}

\begin{solution}{pb:g6:rocks-and-minerals:1}
\textbf{1.} Granite, an \omterm{def:g6:rocks-and-minerals:rock-families}{igneous rock}; its \omterm{def:g6:rocks-and-minerals:mineral}{minerals} are feldspar (pink or
white), quartz (grey, glassy) and mica (black, shiny).

\textbf{2.} Limestone, a \omterm{def:g6:rocks-and-minerals:rock-families}{sedimentary rock}. The shells belonged to animals
that lived in the \omterm{def:g2:water-on-earth:ocean}{sea}; their remains were laid down and hardened into
the \omterm{def:g3:stones-pebbles-sand:rock}{rock}.

\textbf{3.} Slate, a \omterm{def:g6:rocks-and-minerals:rock-families}{metamorphic rock}, formed from shale.

\textbf{4.} Marble, a \omterm{def:g6:rocks-and-minerals:rock-families}{metamorphic rock}, formed from limestone.

\textbf{5.} Quartz: not scratched by steel (about 6.5), it scratches glass
(about 5.5); its \omterm{def:g6:rocks-and-minerals:hardness}{hardness} is 7.

\textbf{6.} Scratched by steel (about 6.5) but not by glass (about 5.5):
between 5.5 and 6.5. It is feldspar (orthoclase, 6).

\textbf{7.} Mica (\omterm{def:g6:rocks-and-minerals:hardness}{hardness} 2 to 3).

\textbf{8.} Yes: the paving is limestone, made of calcite (3), softer than
a copper coin (about 3.5).

\textbf{9.} \omterm{def:g5:water-shapes-the-land:weathering}{Weathering}: the \omterm{def:g4:clouds-rain-snow:precipitation}{rain} slowly wears the surface of the stone
away.

\textbf{10.} The \omterm{def:g1:landscapes-around-us:landform}{mountain} is weathered; its quartz grains are carried by
\omterm{def:g1:landscapes-around-us:water}{rivers} and laid down in layers, at the bottom of a \omterm{def:g1:landscapes-around-us:water}{lake} or of the \omterm{def:g2:water-on-earth:ocean}{sea}; buried
and hardened, they can become a sandstone, a \omterm{def:g6:rocks-and-minerals:rock-families}{sedimentary rock}.

\textbf{11.} All three families are present: igneous (the granite of the
fountain), sedimentary (the limestone of the paving), metamorphic (the
slate of the roof and the marble of the statue).

\textbf{12.} The hardest \omterm{def:g6:rocks-and-minerals:mineral}{mineral} is the quartz of the granite: Mohs
\omterm{def:g6:rocks-and-minerals:hardness}{hardness} \textbf{7}.
\end{solution}
